The objective of this section is to present the variational formulation that will be used for a standard finite-element implementation. 

\subsection{Continuum formulation}
%
%

This section presents the weak forms of the thermo-electro-elastic problem governed by Eqs.~\eqref{eqn:local form conservation of linear momentum}, \eqref{eqn:local form energy our format}, and \eqref{eqn:Gauss law}--\eqref{eqn:Faraday}. 
To this end, we introduce the functional spaces associated with the primary fields 
$\{\vect{\phi}, \varphi, \theta\} \in \{\mathbb{V}^{\vect{\phi}}, \mathbb{V}^{\varphi}, \mathbb{V}^{\theta}\}$, 
together with the corresponding test spaces 
$\{\vect{w}_{\vect{\phi}}, w_{\varphi}, w_{\theta}\} \in \{\mathbb{V}_0^{\vect{\phi}}, \mathbb{V}_0^{\varphi}, \mathbb{V}_0^{\theta}\}$, defined as
%
\begin{equation}\label{eqn:functional spaces}
	\begin{aligned}
		\mathbb{V}^{\vect{\phi}} &=
		\left\{
		\vect{\phi}\in H^1\left(\mathcal{B}_0;\mathbb{R}^3\right)\; \middle|\;
		\vect{\phi} = \bar{\vect{\phi}} \ \text{on}\ \partial_{\vect{\phi}}\mathcal{B}_0
		\right\},
		&
		\mathbb{V}_0^{\vect{\phi}} &=
		\left\{
		\vect{\phi}\in H^1\left(\mathcal{B}_0;\mathbb{R}^3\right)\; \middle|\;
		\vect{\phi} = \vect{0} \ \text{on}\ \partial_{\vect{\phi}}\mathcal{B}_0
		\right\},
		\\
		%
		\mathbb{V}^{\varphi} &=
		\left\{
		\varphi \in H^1\left(\mathcal{B}_0\right)\; \middle|\;
		\varphi = \bar{\varphi} \ \text{on}\ \partial_{\varphi}\mathcal{B}_0
		\right\},
		&
		\mathbb{V}_0^{\varphi} &=
		\left\{
		\varphi \in H^1\left(\mathcal{B}_0\right)\; \middle|\;
		\varphi = 0 \ \text{on}\ \partial_{\varphi}\mathcal{B}_0
		\right\},
		\\
		%
		\mathbb{V}^{\theta} &=
		\left\{
		\theta \in H^1\left(\mathcal{B}_0\right)\; \middle|\;
		\theta = \bar{\theta} \ \text{on}\ \partial_{\theta}\mathcal{B}_0
		\right\},
		&
		\mathbb{V}_0^{\theta} &=
		\left\{
		\theta \in H^1\left(\mathcal{B}_0\right)\; \middle|\;
		\theta = 0 \ \text{on}\ \partial_{\theta}\mathcal{B}_0
		\right\}.
	\end{aligned}
\end{equation}

Following \cite{XX}, the coupled weak forms can be written as
%
\begin{equation}\label{eqn:weak forms for the dynamic formulation}
	\begin{aligned}
		\mathcal{W}_{\vect{\phi}} &=
		\int_{\mathcal{B}_0}\rho_0 \dot{\vect{v}}\cdot\vect{w}_{\vect{\phi}}\,dV 
		+ \int_{\mathcal{B}_0}\vect{P}:\vect{\nabla}_0\vect{w}_{\vect{\phi}}\,dV 
		- \mathcal{W}^{\text{ext}}_{\vect{\phi}} = 0, \\
		%
		\mathcal{W}_{\varphi} &=
		\int_{\mathcal{B}_0}\vect{D}_0\cdot\vect{\nabla}_0 w_{\varphi}\,dV 
		- \mathcal{W}_{\varphi}^{\text{ext}} = 0, \\
		%
		\mathcal{W}_{\theta} &=
		\int_{\mathcal{B}_0}\left(\frac{d}{dt}(\theta\eta) - \dot{\theta}\eta - \mathcal{D}\right) w_{\theta}\,dV 
		- \int_{\mathcal{B}_0}\vect{Q}\cdot\vect{\nabla}_0 w_{\theta}\,dV 
		- \mathcal{W}_{\theta}^{\text{ext}} = 0.
	\end{aligned}
\end{equation}
%
The external mechanical, electrical, and thermal contributions read
%
\begin{equation}
	\begin{aligned}
		\mathcal{W}_{\vect{\phi}}^{\text{ext}} &=
		\int_{\mathcal{B}_0}\vect{f}_0\cdot\vect{w}_{\vect{\phi}}\,dV
		+ \int_{\partial_{\boldsymbol{t}}\mathcal{B}_0}\vect{t}_0\cdot\vect{w}_{\vect{\phi}}\,dA, \\
		%
		\mathcal{W}_{\varphi}^{\text{ext}} &=
		- \int_{\mathcal{B}_0}\rho_0^e w_{\varphi}\,dV
		- \int_{\partial_{\omega}\mathcal{B}_0}\omega_0^e w_{\varphi}\,dA, \\
		%
		\mathcal{W}_{\theta}^{\text{ext}} &=
		\int_{\mathcal{B}_0} R_{\theta} w_{\theta}\,dV
		+ \int_{\partial_{Q}\mathcal{B}_0} Q_{\theta} w_{\theta}\,dA.
	\end{aligned}
\end{equation}


\subsection{Semidiscrete in time formulation}

This section addresses the temporal discretisation of the weak forms in Eq.~\eqref{eqn:weak forms for the dynamic formulation}. 
A second-order accurate trapezoidal rule is employed. 
Consider a sequence of time instants $\{t_1, t_2, \ldots, t_n, t_{n+1}\}$, where $t_{n+1}$ denotes the current time level. 
The kinematic relations are approximated as
%
\begin{equation}\label{eqn:time discretisation kinematics}
	\dot{\vect{v}} \approx \frac{\Delta \vect{v}}{\Delta t},
	\qquad 
	\frac{\Delta \vect{\phi}}{\Delta t} = \vect{v}_{n+1/2},
\end{equation}
%
which yield the equivalent expression
%
\begin{equation*}
	\dot{\vect{v}} \approx 2\frac{\Delta \vect{\phi}}{\Delta t^2} - 2\frac{\vect{v}_n}{\Delta t}.
\end{equation*}

Substitution into Eq.~\eqref{eqn:weak forms for the dynamic formulation} leads to the semi-discrete formulation
%
\begin{equation}
	\label{eqn:weak forms for proposed time integrator}
	\begin{aligned}
		\left(\mathcal{W}_{\vect{\phi}}\right)_{\text{algo}} &=
		\int_{\mathcal{B}_0}\rho_0 
		\left(
		2\frac{\Delta \vect{\phi}}{\Delta t^2} - 2\frac{\vect{v}_n}{\Delta t}
		\right)\cdot\vect{w}_{\vect{\phi}}\,dV 
		+ \int_{\mathcal{B}_0}\vect{P}_{n+1/2}:\vect{\nabla}_0\vect{w}_{\vect{\phi}}\,dV 
		- \left(\mathcal{W}^{\text{ext}}_{\vect{\phi}}\right)_{n+1/2} = 0, \\
		%
		\left(\mathcal{W}_{\varphi}\right)_{\text{algo}} &=
		\int_{\mathcal{B}_0}\vect{D}_{0_{n+1/2}}\cdot\vect{\nabla}_0 w_{\varphi}\,dV 
		+ \int_{\mathcal{B}_0}\rho^{e}_{0_{n+1/2}} w_{\varphi}\,dV 
		- \left(\mathcal{W}_{\varphi}^{\text{ext}}\right)_{n+1/2} = 0, \\
		%
		\left(\mathcal{W}_{\theta}\right)_{\text{algo}} &=
		\int_{\mathcal{B}_0}
		\left(
		\frac{\Delta(\theta\eta)}{\Delta t}
		- \frac{\Delta \theta}{\Delta t}\eta_{n+1/2}
		- \mathcal{D}_{n+1/2}
		\right) w_{\theta}\,dV 
		- \int_{\mathcal{B}_0}\vect{Q}_{n+1/2}\cdot\vect{\nabla}_0 w_{\theta}\,dV 
		- \left(\mathcal{W}_{\theta}^{\text{ext}}\right)_{n+1/2} = 0.
	\end{aligned}
\end{equation}
%
Here, the notation
%
\begin{equation*}
	\Delta(\bullet) = \bullet_{n+1} - \bullet_n,
	\qquad 
	\bullet_{n+1/2} = \frac{1}{2}\left(\bullet_{n+1} + \bullet_n\right)
\end{equation*}
%
is used, where $\bullet_n$ and $\bullet_{n+1}$ denote the evaluation of a generic field $\bullet$ at time levels $t_n$ and $t_{n+1}$, respectively.